\section{Conclusion}\label{sec:conclusion}

This project set out to demonstrate that a fully autonomous search and rescue drone---capable of systematic area coverage, onboard casualty detection, and operator-verified landing---could be built using commodity hardware and open-source software, without reliance on thermal imaging, GPU-class onboard computers, or cloud connectivity. The resulting system achieves this goal, delivering an end-to-end SAR pipeline on a Raspberry Pi~5 companion computer at a total onboard computing cost under \pounds100 and a complete airframe cost of approximately \pounds565.

The principal contributions of this work are fourfold:

\begin{enumerate}
    \item A \textbf{modular finite state machine architecture} with twenty states and thirty-two transitions provides a clear, verifiable structure for the mission lifecycle, from takeoff through lawnmower search to target approach and landing. The strict separation of concerns---with computer vision, path planning, flight control, and configuration each isolated in dedicated modules with narrow interfaces---enabled independent development and testing of each subsystem and allowed the same codebase to run without modification across laptop simulation, bench hardware, and the flight-ready Pi platform.

    \item A \textbf{dual-backend computer vision pipeline} automatically selects between Ultralytics (PyTorch) inference on the development laptop and TensorFlow Lite on the Pi, presenting a single detection interface to the rest of the system. The YOLOv8n model, retrained on 366 mixed synthetic and real-world images at native camera resolution ($1456 \times 1088$), achieves a mean average precision ($\text{mAP}_{50}$) of 0.995 and runs at 206\,ms per frame (4.8\,FPS) on the Pi~5 CPU via XNNPACK---sufficient for the altitude-dependent search speed (8\,m/s at the nominal 35\,m altitude) with overlapping coverage between consecutive frames.

    \item A \textbf{simulation-first development methodology}, progressing through five stages from software-in-the-loop simulation to full autonomous flight readiness, systematically validated each subsystem before integration. This progressive approach caught interface mismatches, timing issues, and parameter errors in simulation that would have been dangerous or difficult to diagnose during live flight.

    \item An \textbf{operator-in-the-loop verification architecture}, implemented through a web-based ground station with MJPEG video streaming and per-target classification controls, ensures that no autonomous landing action is taken without explicit human confirmation, directly addressing the false-positive risk inherent in RGB-only detection from altitude.
\end{enumerate}

Quantitative results from bench testing and video analysis confirm the system's capability. The 18-waypoint lawnmower pattern covers the designated search polygon in approximately 1.1\,minutes at the nominal 8\,m/s search speed. GPS-based target localisation, combining pixel-to-ground projection with spatial clustering and Kalman filtering, achieves a circular error probable (CEP$_{50}$) of 2.3\,m from DJI flight video analysis---within the accuracy required for guiding ground teams to a casualty's location. Lens distortion correction, calibrated to an RMS reprojection error of 0.399\,pixels, adds only 1.5\,ms per frame to the processing pipeline.

Several limitations should be acknowledged. The training dataset, while achieving high validation accuracy, comprises only 366 images and has not been evaluated against the full diversity of outdoor conditions---varied lighting, weather, vegetation, and casualty poses---that an operational system would encounter. Inference remains single-threaded, leaving potential performance gains from pipelined capture and processing unexploited. The system relies exclusively on RGB imagery, which limits detection to daylight conditions with adequate visual contrast between the casualty and the background; thermal imaging would extend capability to low-visibility and night-time scenarios but was excluded to maintain the low-cost design constraint.

Despite these limitations, the system demonstrates that effective autonomous SAR capability is achievable with consumer-grade hardware and open-source software. The modular architecture, progressive testing methodology, and operator-in-the-loop design provide a foundation that can be extended with faster inference backends, expanded training data, thermal sensor fusion, and multi-drone coordination, as discussed in Section~\ref{sec:future}. By lowering the cost and complexity barriers to autonomous aerial search, this work contributes toward making life-saving UAV technology accessible to the volunteer and community rescue organisations that need it most.
